\documentclass[12pt,a4paper]{article}

% --- Packages ---
\usepackage[T1]{fontenc}
\usepackage[utf8]{inputenc}
\usepackage[indonesian]{babel}
\usepackage{geometry}
\usepackage{amsmath,amssymb,amsfonts}
\usepackage{graphicx}
\usepackage{xcolor}
\usepackage{booktabs}
\usepackage{tabularx}
\usepackage{array}
\usepackage{multirow}
\usepackage{hyperref}
\usepackage{fancyhdr}
\usepackage{titlesec}
\usepackage{tcolorbox}
\usepackage{enumitem}
\usepackage{listings}
\usepackage{courier}
\usepackage{mathtools}
\usepackage{caption}
\usepackage{tikz}
\usepackage{pgfplots}
\pgfplotsset{compat=1.18}
\usetikzlibrary{arrows.meta, positioning, shapes.geometric, backgrounds}

% --- Page Geometry ---
\geometry{top=2.5cm, bottom=2.5cm, left=3cm, right=2.5cm, headheight=15pt}

% --- Colors ---
\definecolor{primaryblue}{RGB}{37,99,235}
\definecolor{accentcyan}{RGB}{6,182,212}
\definecolor{successgreen}{RGB}{52,211,153}
\definecolor{darkgray}{RGB}{60,60,60}
\definecolor{sectioncolor}{RGB}{30,58,138}
\definecolor{boxblue}{RGB}{219,234,254}
\definecolor{boxblueborder}{RGB}{147,197,253}
\definecolor{codegray}{RGB}{248,248,248}

% --- Hyperref ---
\hypersetup{
  colorlinks=true,
  linkcolor=primaryblue,
  urlcolor=accentcyan,
  pdftitle={Rangkuman TA: Simulator Platooning V2V Berbasis 5G},
  pdfauthor={Mahendra Aryaputra Fitrianto},
}

% --- Section Formatting ---
\titleformat{\section}{\large\bfseries\color{sectioncolor}}{\thesection.}{0.6em}{}[\titlerule]
\titleformat{\subsection}{\normalsize\bfseries\color{primaryblue}}{\thesubsection.}{0.5em}{}
\titleformat{\subsubsection}{\normalsize\bfseries\color{darkgray}}{\thesubsubsection.}{0.5em}{}

% --- Header & Footer ---
\pagestyle{fancy}
\fancyhf{}
\fancyhead[L]{\small\color{gray}Tugas Akhir -- Simulator Platooning V2V Berbasis 5G}
\fancyhead[R]{\small\color{gray}S1 Teknik Telekomunikasi, Telkom University}
\fancyfoot[C]{\small\thepage}
\renewcommand{\headrulewidth}{0.4pt}

% --- Code Listing ---
\lstdefinestyle{codeStyle}{
  backgroundcolor=\color{codegray},
  basicstyle=\ttfamily\footnotesize,
  keywordstyle=\color{primaryblue}\bfseries,
  commentstyle=\color{gray}\itshape,
  stringstyle=\color{successgreen},
  frame=single,
  framesep=6pt,
  rulecolor=\color{boxblueborder},
  breaklines=true,
  numbers=left,
  numberstyle=\tiny\color{gray},
  numbersep=8pt,
  xleftmargin=18pt,
  tabsize=2,
}

% --- tcolorbox Styles ---
\tcbuselibrary{skins,breakable}
\newtcolorbox{formulabox}[1]{
  colback=boxblue, colframe=primaryblue,
  title={\bfseries #1},
  fonttitle=\small\bfseries\color{white},
  coltitle=white, boxrule=1pt, arc=4pt,
  left=6pt, right=6pt, top=4pt, bottom=4pt, breakable
}
\newtcolorbox{infobox}[1]{
  colback=yellow!10, colframe=orange!70!black,
  title={\bfseries #1},
  fonttitle=\small\bfseries,
  boxrule=0.8pt, arc=4pt,
  left=6pt, right=6pt, top=4pt, bottom=4pt, breakable
}
\newtcolorbox{successbox}[1]{
  colback=green!8, colframe=successgreen,
  title={\bfseries #1},
  fonttitle=\small\bfseries,
  boxrule=0.8pt, arc=4pt,
  left=6pt, right=6pt, top=4pt, bottom=4pt, breakable
}

% ============================================================
\begin{document}
% ============================================================

% --- Cover Page ---
\begin{titlepage}
  \centering
  \vspace*{1cm}
  {\Large\bfseries\color{sectioncolor} RANGKUMAN TUGAS AKHIR\\[0.5em]}
  \rule{\linewidth}{2pt}\\[0.8em]
  {\Huge\bfseries\color{primaryblue}
    Perancangan Sistem Simulasi\\
    Platooning \textit{Vehicle-to-Vehicle} (V2V)\\
    Berbasis 5G\\[0.5em]}
  \rule{\linewidth}{2pt}\\[1.5em]

  {\large
    Platform Simulasi 2D Interaktif \textit{Client-Server}\\
    Terintegrasi Emulator Jaringan 5G URLLC\\
    dan Algoritma Kontrol CACC/ACC
  }\\[2.5em]

  \begin{tabular}{rl}
    \textbf{Penulis:}     & Mahendra Aryaputra Fitrianto \\
                          & Muhammad Abduh \\
                          & Ahmad Zulfikar \\[0.8em]
    \textbf{Pembimbing:}  & Linda Meylani S.T., M.T. \\
                          & Ir. Uke Kurniawan Usman, M.T. \\[0.8em]
    \textbf{Program Studi:} & S1 Teknik Telekomunikasi \\
    \textbf{Fakultas:}    & Teknik Elektro \\
    \textbf{Universitas:} & Telkom University \\[0.8em]
    \textbf{Tahun:}       & 2026 \\
  \end{tabular}

  \vfill
  {\small\color{gray}
    Dokumen ini adalah rangkuman akademik komprehensif yang mencakup seluruh konsep,\\
    arsitektur, formula matematis, dan implementasi teknis dari Tugas Akhir di atas.
  }
\end{titlepage}

\tableofcontents
\newpage

% ============================================================
\section{Abstrak}
% ============================================================

Sistem \textit{vehicular platooning} meningkatkan efisiensi lalu lintas melalui koordinasi
otomatis, namun rentan terhadap fluktuasi \textit{latency} dan \textit{packet loss}.
Guna menghindari bahaya pengujian fisik, penelitian ini merancang platform simulasi 2D
interaktif \textit{Client-Server} terintegrasi emulator 5G URLLC dan logika CACC.
Platform ini memvalidasi keamanan formasi di bawah gangguan jaringan dinamis dan menampilkan
analisis stabilitas secara \textit{real-time}.

\begin{successbox}{Kontribusi Utama Penelitian}
\begin{enumerate}[noitemsep]
  \item Platform simulasi 2D \textit{client-server} berhasil memvalidasi algoritma kontrol
        CACC secara aman, presisi, dan terukur tanpa risiko fisik.
  \item Integrasi emulator 5G memetakan pengaruh degradasi jaringan (latensi \&
        \textit{packet loss}) terhadap stabilitas barisan (\textit{string stability}).
  \item Kontrol PD \textit{feedforward} terbukti efektif meminimalkan \textit{spacing error}
        platoon di bawah gangguan jaringan nirkabel dinamis.
\end{enumerate}
\end{successbox}

% ============================================================
\section{Latar Belakang dan Motivasi Penelitian}
% ============================================================

\subsection{Vehicular Platooning dan Intelligent Transportation System (ITS)}

Perkembangan kendaraan otonom dan komunikasi \textit{Vehicle-to-Vehicle} (V2V) merupakan
fokus utama \textit{Intelligent Transportation System} (ITS) modern. \textit{Platooning}
adalah formasi barisan kendaraan otomatis yang melaju bersama dengan jarak rapat seperti
rangkaian kereta di jalan raya.

\textbf{Manfaat utama platooning:}
\begin{itemize}[noitemsep]
  \item Penghematan bahan bakar 10--15\% akibat berkurangnya hambatan aerodinamik.
  \item Peningkatan kapasitas jalan raya dan pengurangan kemacetan.
  \item Pencegahan kelelahan pengemudi pada perjalanan jauh.
  \item Potensi reduksi emisi gas buang secara signifikan.
\end{itemize}

\subsection{Keterbatasan Adaptive Cruise Control (ACC) Konvensional}

\textit{Adaptive Cruise Control} (ACC) konvensional hanya mengandalkan sensor lokal seperti
radar atau LiDAR. Karena kendaraan hanya dapat melihat kendaraan tepat di depannya, terjadi
fenomena yang disebut \textbf{\textit{string instability}}: gangguan pengereman di depan
diamplifikasi (diperbesar) ke kendaraan di belakang, sehingga pengereman mendadak di bagian
depan konvoi berisiko menyebabkan tabrakan beruntun di bagian belakang.

\subsection{Teknologi 5G URLLC sebagai Enabler CACC}

Teknologi 5G \textit{Ultra-Reliable Low-Latency Communication} (URLLC) menjanjikan:
\begin{itemize}[noitemsep]
  \item Latensi \textit{end-to-end} di bawah 1 ms (teoritis) hingga 10--20 ms (praktis).
  \item Keandalan paket data di atas 99{,}999\%.
  \item Kemampuan transmisi data multi-kendaraan secara simultan dalam frekuensi tinggi.
\end{itemize}

Dengan memanfaatkan 5G V2X, kendaraan pengikut tidak hanya mengandalkan radar, tetapi
juga menerima informasi instan mengenai percepatan, kecepatan, dan posisi kendaraan lain
secara \textit{real-time}.

\subsection{Rumusan Masalah}

Bagaimana komunikasi V2V berbasis 5G memengaruhi kestabilan sistem kontrol
\textit{platooning}, serta bagaimana fungsi-fungsi di dalamnya dapat diintegrasikan agar
tetap sinkron meskipun jaringan mengalami degradasi performa (latensi dan
\textit{packet loss})?

\subsection{Tujuan Penelitian}

Mengembangkan platform pengujian (simulasi) yang mampu memvalidasi algoritma
\textit{Cooperative Adaptive Cruise Control} (CACC) dan protokol komunikasi V2V secara
presisi, aman, dan terukur \textbf{tanpa adanya risiko kerusakan fisik}.

% ============================================================
\section{Arsitektur Sistem}
% ============================================================

\subsection{Desain Arsitektur Client-Server}

Sistem mengadopsi desain \textbf{arsitektur Client-Server} berbasis web yang memisahkan
logika komputasi (backend) dari lapisan visualisasi (frontend).

\begin{table}[ht]
\centering
\caption{Lapisan Arsitektur Sistem}
\begin{tabularx}{\linewidth}{lX}
\toprule
\textbf{Layer} & \textbf{Deskripsi} \\
\midrule
\textbf{Client (Frontend)} &
  React 18 + Vite + TypeScript. Visualisasi Canvas 2D pergerakan kendaraan
  real-time, panel konfigurasi parameter, Telemetry HUD. \\[0.5em]
\textbf{5G Network Emulator} &
  Emulasi latensi (buffer antrean) dan packet loss (probabilistic drop).
  Dijalankan di sisi backend sebagai bagian dari loop fisika. \\[0.5em]
\textbf{Server (Backend)} &
  Node.js + Express + TypeScript. Physics engine CACC/ACC beroperasi pada
  20 Hz. Mengelola state kendaraan, emulasi jaringan, dan API Socket.IO. \\[0.5em]
\textbf{Database \& Session Log} &
  Penyimpanan riwayat sesi (JSON), telemetri berseri, dan hasil analisis
  stabilitas yang dapat diekspor. \\
\bottomrule
\end{tabularx}
\end{table}

Komunikasi antara frontend dan backend menggunakan \textbf{WebSocket via Socket.IO},
memungkinkan transmisi state simulasi dua arah dengan latensi rendah pada frekuensi
$\approx$20 Hz.

\subsection{Stack Teknologi}

\begin{table}[ht]
\centering
\caption{Stack Teknologi Implementasi}
\begin{tabularx}{\linewidth}{lXX}
\toprule
\textbf{Layer} & \textbf{Teknologi} & \textbf{Fungsi} \\
\midrule
Frontend & React 18 + Vite + TypeScript & SPA, rendering Canvas 2D, HUD Telemetry \\
Backend  & Node.js + Express + TypeScript & Physics engine, CACC/ACC loop, REST API \\
Real-Time Comm. & Socket.IO (WebSocket) & Transmisi state simulasi 20 Hz \\
Visualisasi & HTML5 Canvas 2D API & Render kendaraan, lajur, RSU \\
Deployment & PM2 (VPS DigitalOcean) & Process manager produksi \\
Reverse Proxy & Nginx + Let's Encrypt SSL & HTTPS \& domain routing \\
\bottomrule
\end{tabularx}
\end{table}

\subsection{Alur Data Real-Time}

\begin{enumerate}
  \item Backend menjalankan \textit{physics loop} konstan pada $\approx$20 Hz (setiap 50 ms).
  \item Setiap \textit{tick}, backend menghitung posisi, kecepatan, dan akselerasi setiap
        kendaraan berdasarkan model CACC/ACC.
  \item State lengkap dikirim ke semua klien melalui event Socket.IO (\texttt{sim:state}).
  \item Frontend menerima data dan merender ulang canvas serta memperbarui panel telemetri.
  \item Pengguna mengubah parameter simulasi melalui slider dan dikirim balik ke backend via
        event \texttt{sim:updateParams}.
\end{enumerate}

% ============================================================
\section{Model Kendaraan dan Kinematika}
% ============================================================

\subsection{State Kendaraan}

Setiap kendaraan dimodelkan sebagai objek \texttt{VehicleState}:

\begin{table}[ht]
\centering
\caption{Atribut State Kendaraan (\texttt{VehicleState})}
\begin{tabularx}{\linewidth}{llX}
\toprule
\textbf{Atribut} & \textbf{Tipe} & \textbf{Keterangan} \\
\midrule
\texttt{x}       & number (m) & Posisi longitudinal di sepanjang jalan (meter) \\
\texttt{y}       & number (int) & Indeks lajur platoon (0 = Platoon A, 1 = B, dst.) \\
\texttt{wy}      & number & Posisi lateral kontinu dalam satuan lajur \\
\texttt{heading} & number (rad) & Sudut heading kendaraan; 0 = lurus ke depan \\
\texttt{speed}   & number (m/s) & Kecepatan longitudinal \\
\texttt{accel}   & number (m/s$^{2}$) & Percepatan/deselerasi saat ini \\
\texttt{brake}   & boolean & Indikator apakah kendaraan sedang mengerem \\
\texttt{crashed} & boolean & Tanda tabrakan; kendaraan dibekukan jika \texttt{true} \\
\bottomrule
\end{tabularx}
\end{table}

\subsection{Konstanta Fisika}

\begin{table}[ht]
\centering
\caption{Konstanta Fisika Simulasi}
\begin{tabular}{lll}
\toprule
\textbf{Konstanta} & \textbf{Nilai} & \textbf{Keterangan} \\
\midrule
$a_{\max}$ & 3.5 m/s$^{2}$ & Batas akselerasi mesin (throttle penuh) \\
$b_{\max}$ & 6.0 m/s$^{2}$ & Batas deselerasi normal \\
$b_{\text{emrg}}$ & 9.0 m/s$^{2}$ & Deselerasi darurat saat jarak kritis \\
$v_{\max}$ & 42.0 m/s ($\approx$150 km/h) & Kecepatan maksimum kendaraan \\
$w_{\text{lane}}$ & 3.5 m & Lebar lajur jalan nyata \\
$\Delta t$ & 0.05 s & Interval waktu simulasi (20 Hz) \\
\bottomrule
\end{tabular}
\end{table}

\subsection{Pembaruan Posisi Kendaraan (Integrasi Euler)}

Sistem menggunakan metode integrasi \textbf{Euler Forward}:

\begin{formulabox}{Kinematika Longitudinal}
\begin{align}
  v(t{+}\Delta t) &= \operatorname{clamp}\!\bigl(v(t) + a_{\text{cmd}} \cdot \Delta t,\; 0,\; v_{\max}\bigr) \label{eq:v_update}\\[4pt]
  x(t{+}\Delta t) &= x(t) + v(t{+}\Delta t) \cdot \cos(\theta) \cdot \Delta t \label{eq:x_update}
\end{align}
di mana $\theta$ adalah sudut \textit{heading} kendaraan dan $\Delta t = 0.05$ s.
\end{formulabox}

\subsection{Model Kinematika Lateral (Lane Change)}

Saat terjadi pergantian lajur, sistem menggunakan \textbf{trajektori sinusoidal halus}:

\begin{formulabox}{Trajektori Sinusoidal Pergantian Lajur}
\begin{align}
  p(t) &= \tfrac{1}{2} - \tfrac{1}{2}\cos\!\!\left(\pi \cdot \frac{t_{\text{elapsed}}}{T_{\text{dur}}}\right) \label{eq:sinusoidal}\\[4pt]
  y_w(t) &= y_{\text{start}} + (y_{\text{target}} - y_{\text{start}}) \cdot p(t) \label{eq:lane_pos}\\[4pt]
  \theta(t) &= \arctan\!\!\left(\frac{\Delta y_w \cdot w_{\text{lane}}}{v \cdot \Delta t}\right) \label{eq:heading}
\end{align}
di mana $p(t) \in [0,1]$ adalah progress maneuver, $T_{\text{dur}}$ = durasi total maneuver.
\end{formulabox}

Untuk gerakan lateral normal (proportional steering fallback):
\begin{equation}
  \theta_{\text{des}} = \operatorname{clamp}\!\bigl(K_{\text{steer}} \cdot (y_{\text{target}} - y_w),\; -\theta_{\max},\; \theta_{\max}\bigr), \quad K_{\text{steer}} = 1.5,\; \theta_{\max} = 0.12\text{ rad}
\end{equation}

% ============================================================
\section{Algoritma Kontrol CACC dan ACC}
% ============================================================

\subsection{Desired Gap Model (Constant Time Headway Policy)}

Kedua kontroller (CACC dan ACC) menggunakan model \textbf{CTHP} sebagai acuan jarak aman:

\begin{formulabox}{Desired Gap -- Jarak Aman Target}
\begin{equation}
  d_{\text{desired}} = s_0 + h \cdot v_f \label{eq:d_desired}
\end{equation}
\begin{itemize}[noitemsep]
  \item $s_0$ = standstill distance (jarak minimum saat berhenti) [m], rentang: 4--20 m
  \item $h$ = time headway (waktu jeda aman) [s], rentang: 0.5--3.0 s
  \item $v_f$ = kecepatan kendaraan pengikut (\textit{follower}) [m/s]
\end{itemize}
\end{formulabox}

\subsection{Spacing Error dan Kecepatan Relatif}

\begin{formulabox}{Spacing Error (Eror Jarak) dan Kecepatan Relatif}
\begin{align}
  d_{\text{actual}} &= x_{\text{pred}} - x_f \label{eq:d_actual}\\[4pt]
  e &= d_{\text{actual}} - d_{\text{desired}} \label{eq:spacing_err}\\[4pt]
  v_{\text{rel}} &= v_{\text{pred}} - v_f \label{eq:v_rel}
\end{align}
$e > 0$: kendaraan terlalu jauh (boleh akselerasi);\quad $e < 0$: terlalu dekat (harus rem).
\end{formulabox}

\subsection{Algoritma CACC (Cooperative Adaptive Cruise Control)}

CACC menggabungkan umpan balik lokal dengan \textbf{umpan maju nirkabel (feedforward)}:

\begin{formulabox}{Perintah Akselerasi CACC -- PD + Feedforward}
\begin{equation}
  \boxed{a_{\text{CACC}} = K_P \cdot e + K_D \cdot v_{\text{rel}} + \alpha_{\text{ff}} \cdot a_{\text{ff}}}
  \label{eq:cacc}
\end{equation}
\begin{center}
\begin{tabular}{lll}
  $K_P = 0.45$ & Gain proporsional & (mengoreksi eror jarak) \\
  $K_D = 0.35$ & Gain derivatif & (mengoreksi perbedaan kecepatan) \\
  $\alpha_{\text{ff}} = 0.6$ & Bobot feedforward & (data akselerasi V2V nirkabel) \\
\end{tabular}
\end{center}
Output di-clamp ke batas fisik: $a_{\text{cmd}} = \operatorname{clamp}(a_{\text{CACC}}, -b_{\max}, a_{\max})$
\end{formulabox}

\subsection{Topologi Komunikasi V2V}

Sumber data $a_{\text{ff}}$ bergantung pada topologi yang dipilih:

\begin{table}[ht]
\centering
\caption{Perbandingan Topologi Komunikasi V2V}
\begin{tabularx}{\linewidth}{lXX}
\toprule
\textbf{Topologi} & \textbf{Sumber $a_{\text{ff}}$} & \textbf{Karakteristik} \\
\midrule
\textbf{Predecessor Following (PF)} &
  $a_{\text{predecessor}}$ (mobil langsung di depan) &
  Informasi berantai; rentan akumulasi delay pada platoon panjang \\[0.5em]
\textbf{Leader-to-All (L2A)} &
  $a_{\text{leader}}$ (mobil terdepan) &
  Semua pengikut bereaksi seketika; latensi tinggi jadi kelemahan \\[0.5em]
\textbf{Hybrid (Default)} &
  $a_{\text{leader}}$ untuk feedforward; jarak acuan ke predecessor &
  Kombinasi terbaik: respons cepat + keamanan jarak lokal \\
\bottomrule
\end{tabularx}
\end{table}

\subsection{ACC Fallback (Radar-only Mode)}

Ketika $\text{PLR} \geq 15\%$, sistem otomatis beralih ke mode ACC konvensional:

\begin{formulabox}{Perintah Akselerasi ACC Fallback}
\begin{align}
  h_{\text{safe}} &= \max(2.0,\; h \times 1.6) \label{eq:acc_h}\\[4pt]
  d_{\text{des,ACC}} &= s_0 + h_{\text{safe}} \cdot v_f \label{eq:acc_d}\\[4pt]
  a_{\text{ACC}} &= 0.32 \cdot e_{\text{ACC}} + 0.18 \cdot v_{\text{rel}} \cdot 0.4 \label{eq:acc_a}
\end{align}
Tanpa data prediktif V2V, headway diperlebar $1.6\times$ (min. 2.0 s) untuk mengakomodasi
batas \textit{stopping distance} radar. Gain diturunkan agar rem tidak kejut.
\end{formulabox}

\begin{infobox}{Perbandingan CACC vs ACC Fallback}
\begin{center}
\begin{tabular}{lcc}
\toprule
\textbf{Parameter} & \textbf{CACC} & \textbf{ACC Fallback} \\
\midrule
Sumber Informasi & Radar + V2V Nirkabel & Radar Lokal Saja \\
Gain $K_P$ / $K_D$ & 0.45 / 0.35 & 0.32 / 0.072 \\
Umpan maju $\alpha_{\text{ff}}$ & 0.6 & 0 \\
Time Headway & $h$ (konfigurabel) & $\max(2.0,\; 1.6h)$ \\
Aktivasi & Default & PLR $\geq 15\%$ \\
\bottomrule
\end{tabular}
\end{center}
\end{infobox}

\subsection{Kontrol Leader (Kendaraan Pemimpin)}

\begin{equation}
  a_{\text{leader}} = \tau \cdot a_{\max} - \beta \cdot b_{\max}, \qquad
  v_{\text{leader}} \leftarrow \operatorname{clamp}(v + a_{\text{leader}} \cdot \Delta t,\; 0,\; v_{\max})
\end{equation}
di mana $\tau \in [0,1]$ adalah throttle dan $\beta \in [0,1]$ adalah input rem operator.

% ============================================================
\section{Emulasi Jaringan 5G}
% ============================================================

\subsection{Arsitektur Emulator Jaringan}

Emulator jaringan 5G menggunakan dua mekanisme utama:
\begin{enumerate}
  \item \textbf{Antrean Penundaan (Latency Buffer Queue)}: Paket disimpan dengan stempel
        waktu pengantaran $t_{\text{deliver}} = t_{\text{now}} + L_{\text{ms}}$.
  \item \textbf{Pembuangan Probabilistik (Probabilistic Drop)}: Jika $r < \text{PLR}$,
        di mana $r \sim \mathcal{U}(0,100)$, paket dibuang.
\end{enumerate}

\begin{formulabox}{Model Emulasi Paket V2V}
\begin{align}
  t_{\text{deliver}} &= t_{\text{now}} + L_{\text{ms}} \label{eq:latency}\\[4pt]
  \text{packet lost} &= \begin{cases}
    \text{true}  & \text{jika } r < \text{PLR},\quad r \sim \mathcal{U}(0,100) \\
    \text{false} & \text{lainnya}
  \end{cases} \label{eq:plr}
\end{align}
\end{formulabox}

\subsection{Parameter Jaringan yang Dapat Dikonfigurasi}

\begin{table}[ht]
\centering
\caption{Parameter Jaringan Simulasi}
\begin{tabularx}{\linewidth}{lllX}
\toprule
\textbf{Parameter} & \textbf{Simbol} & \textbf{Rentang} & \textbf{Keterangan} \\
\midrule
One-Way Latency & $L$ & 5--500 ms & Keterlambatan transmisi paket satu arah \\
Packet Loss Rate & PLR & 0--100\% & Persentase paket yang hilang di jaringan \\
Channel Bandwidth & $B$ & 5 MHz--5 GHz & Lebar pita saluran 5G NR \\
V2V Topology & --- & PF, L2A, Hybrid & Pola komunikasi antar kendaraan \\
Dynamic Path Loss & --- & On/Off & Model 3GPP berbasis jarak RSU \\
\bottomrule
\end{tabularx}
\end{table}

\subsection{Model Dynamic Path Loss (Standar 3GPP)}

Ketika \textit{Dynamic Path Loss} aktif, PLR dihitung berdasarkan jarak ke RSU terdekat:

\begin{formulabox}{Model Path Loss 3GPP (Terinspirasi)}
\begin{equation}
  \text{PLR}_{\text{dyn}}(d) = \begin{cases}
    0\%   & d \leq 50\text{ m (Zona Dekat)} \\
    80\% \cdot \bigl(1 - e^{-k(d-50)}\bigr) & 50 < d < 300\text{ m (Zona Menengah)} \\
    80\%  & d \geq 300\text{ m (Zona Jauh)}
  \end{cases}
  \label{eq:path_loss}
\end{equation}
RSU ditempatkan berkala setiap 500 m, jangkauan efektif 300 m.
\end{formulabox}

\subsection{Utilisasi Saluran dan E2E Delay}

\begin{formulabox}{Utilisasi Kanal dan End-to-End Delay}
\begin{align}
  U &= \min\!\left(100,\;\frac{N_v \times 0.05}{B_{\text{MHz}}} \times 100\right) \label{eq:util}\\[6pt]
  T_{\text{E2E}} &= L_{\text{ms}} \times \max(1,\; N_{\text{platoon}} - 1) \label{eq:e2e}
\end{align}
$U$ = utilisasi kanal (\%), $N_v$ = jumlah kendaraan, $B_{\text{MHz}}$ = bandwidth dalam MHz.\\
$T_{\text{E2E}}$ = total delay multi-hop dari Leader ke Follower terakhir.
\end{formulabox}

\subsection{Deteksi Tabrakan}

\begin{formulabox}{Deteksi Tabrakan (Bounding Box)}
\begin{equation}
  \text{collision} = |x_A - x_B| < 4.5\text{ m} \;\wedge\; |y_A - y_B| < 2.0\text{ m}
  \label{eq:collision}
\end{equation}
Saat tabrakan terdeteksi: kendaraan dibekukan, simulasi dihentikan otomatis,
sesi disimpan untuk analisis.
\end{formulabox}

% ============================================================
\section{Metrik Evaluasi Kinerja (KPI)}
% ============================================================

\subsection{String Stability Index (SSI)}

\begin{formulabox}{String Stability Index (SSI)}
\begin{equation}
  \boxed{SSI = \max\!\left(0,\; 1 - \frac{|e_{\text{total}}|}{20}\right)}
  \label{eq:ssi}
\end{equation}
\begin{itemize}[noitemsep]
  \item $SSI = 1.0$: Platoon sempurna stabil, jarak aman dijaga presisi.
  \item $0.8 \leq SSI < 1.0$: Sedikit ketidakstabilan; masih dapat diterima.
  \item $SSI < 0.8$: Ketidakstabilan serius (\textit{string instability}); risiko tabrakan tinggi.
  \item $SSI = 0$: Eror jarak mencapai atau melebihi 20 meter.
\end{itemize}
\end{formulabox}

Jika pengereman di Leader menghasilkan eror jarak yang semakin besar saat merambat ke
belakang, SSI akan menurun. Kontrol CACC yang efektif mempertahankan SSI mendekati 1.0.

\subsection{Ringkasan KPI Simulasi}

\begin{table}[ht]
\centering
\caption{Key Performance Indicators (KPI) Simulasi}
\begin{tabularx}{\linewidth}{lXll}
\toprule
\textbf{KPI} & \textbf{Definisi} & \textbf{Satuan} & \textbf{Ideal} \\
\midrule
SSI & Indeks kestabilan formasi barisan & -- & $\to 1.0$ \\
Avg. Spacing Error & Rata-rata eror jarak antar kendaraan & m & $\to 0$ \\
Max Spacing Error & Eror jarak maksimum yang pernah terjadi & m & $\to 0$ \\
E2E Delay & Keterlambatan total Leader ke Follower terakhir & ms & Rendah \\
Avg. Latency & Rata-rata latensi jaringan per tick & ms & $<20$ ms \\
Avg. PLR & Rata-rata paket yang hilang & \% & $<5\%$ \\
ACC Fallback \% & Waktu simulasi dalam mode ACC & \% & $\to 0\%$ \\
Avg. Update Hz & Frekuensi rata-rata pembaruan simulasi & Hz & $\approx 20$ \\
Collision Count & Jumlah tabrakan yang terjadi & kejadian & 0 \\
\bottomrule
\end{tabularx}
\end{table}

% ============================================================
\section{Skenario Eksperimen}
% ============================================================

Platform mendukung tiga jenis gangguan yang dapat dipicu secara interaktif:

\begin{table}[ht]
\centering
\caption{Skenario Gangguan Interaktif}
\begin{tabularx}{\linewidth}{llX}
\toprule
\textbf{Hotkey} & \textbf{Nama Skenario} & \textbf{Deskripsi} \\
\midrule
\texttt{[1]} & Latency Spike & Lonjakan latensi tiba-tiba ke jaringan V2V \\
\texttt{[2]} & Packet Drop & Kehilangan paket masif secara seketika \\
\texttt{[3]} & Human Braking & Kendaraan Leader melakukan pengereman mendadak \\
\bottomrule
\end{tabularx}
\end{table}

\subsection{Skenario 1: CACC dalam Kondisi Ideal (5G Sempurna)}

\textbf{Konfigurasi}: $L = 5$ ms, PLR = 0\%, $B = 100$ MHz, Topology = Hybrid.

\textbf{Hasil yang Diharapkan}: Platoon merespons pengereman secara simultan dan
terkoordinasi. SSI tetap mendekati 1.0. Barisan melambat rapi tanpa eror jarak signifikan.

\subsection{Skenario 2: Degradasi Jaringan Dinamis (RSU Path Loss)}

\textbf{Konfigurasi}: Aktifkan \textit{Dynamic Path Loss (3GPP)}.

\textbf{Observasi}: Saat kendaraan menjauhi tiang RSU ($d > 250$ m), PLR melonjak.
Indikator V2V berubah \textit{Connected} $\to$ \textit{Degraded} $\to$
\textit{Disconnected}. ACC Fallback aktif dan jarak aman melebar.

\textbf{Pemulihan}: Mendekati RSU berikutnya, CACC pulih otomatis dan jarak kembali normal.

\subsection{Skenario 3: Kegagalan Formasi (String Instability)}

\textbf{Konfigurasi}: $L = 20$ ms, PLR = 14\%, Topology = PF.

\textbf{Hasil}: Efek domino amplifikasi eror berujung tabrakan di kendaraan belakang.
Layar berkedip merah dan \textit{Crash Modal} ditampilkan. SSI turun mendekati 0.

% ============================================================
\section{Implementasi Teknis}
% ============================================================

\subsection{Struktur Proyek}

\begin{lstlisting}[style=codeStyle, caption=Struktur Folder Proyek]
TA/
  backend/src/sim/
    cacc.ts           -- Algoritma CACC & ACC Fallback
    networkEmulator.ts -- 5G latency queue & packet drop
    physics.ts        -- Kinematika & pembaruan posisi kendaraan
    sessionManager.ts -- Log telemetri & penyimpanan histori
    types.ts          -- Tipe data TypeScript (state, params)
    validation.ts     -- Sanitasi & validasi input parameter
  backend/src/
    server.ts         -- Entry point: Socket.IO, REST API, loop fisika
  frontend/src/
    components/       -- ControlPanel, SimulationCanvas, TelemetryPanel
    pages/            -- SimulationPage, AnalysisPage
    hooks/            -- useSimulationSocket (Socket.IO management)
  ecosystem.config.js  -- Konfigurasi PM2 (produksi VPS)
  run.bat             -- Launcher lokal otomatis (Windows)
\end{lstlisting}

\subsection{Implementasi Inti CACC}

\begin{lstlisting}[style=codeStyle, caption=Inti Algoritma CACC (cacc.ts)]
const KP = 0.45   // spacing error proportional gain
const KD = 0.35   // relative speed derivative gain

function computeCaccAcceleration(input):
  desiredGap   = input.standstillDistance
               + input.timeHeadway * input.followerSpeed
  actualGap    = input.predecessorX - input.followerX
  spacingError = actualGap - desiredGap
  relativeSpeed = input.predecessorSpeed - input.followerSpeed

  // Topology-based feedforward
  switch (input.topology):
    case 'PF':     feedforward = input.predecessorAccel
    case 'L2A':    feedforward = input.leaderAccel
    case 'Hybrid': feedforward = input.leaderAccel

  accelCmd = KP * spacingError
           + KD * relativeSpeed
           + feedforward * 0.6

  return { accelCmd, spacingError }
\end{lstlisting}

\subsection{Implementasi Emulator Jaringan}

\begin{lstlisting}[style=codeStyle, caption=Emulator Jaringan 5G (networkEmulator.ts)]
class NetworkEmulator:
  queue = []          // antrean paket dengan stempel waktu
  lastGoodPacket = null

  push(packet, latencyMs, packetLossPercent):
    isLost = (Math.random() * 100) < packetLossPercent
    if isLost: return     // paket dibuang (drop)
    queue.push({ deliverAt: now() + latencyMs, packet })

  receive():
    idx = queue.findIndex(item => item.deliverAt <= now())
    if idx < 0: return lastGoodPacket  // gunakan data lama
    item = queue.splice(idx, 1)
    lastGoodPacket = item.packet
    return item.packet
\end{lstlisting}

\subsection{Infrastruktur Deployment}

\begin{table}[ht]
\centering
\caption{Konfigurasi Infrastruktur Deployment}
\begin{tabularx}{\linewidth}{lX}
\toprule
\textbf{Komponen} & \textbf{Konfigurasi} \\
\midrule
VPS Provider & DigitalOcean (188.166.178.227) \\
OS & Ubuntu 22.04 LTS \\
Process Manager & PM2 (\texttt{v2v-platooning-simulation}) \\
Reverse Proxy & Nginx dengan SSL/TLS (Let's Encrypt) \\
Domain Publik & \texttt{https://v2v-platooning-simulator.cloud} \\
Port Backend & 4000 (internal), 443 HTTPS (publik via Nginx) \\
Build Pipeline & \texttt{npm run build}: TypeScript $\to$ JavaScript dist \\
Deploy Method & SFTP upload \texttt{frontend/dist} + \texttt{backend/dist}, restart PM2 \\
\bottomrule
\end{tabularx}
\end{table}

% ============================================================
\section{Rekap Seluruh Formula Matematis}
% ============================================================

\begin{table}[ht]
\centering
\caption{Daftar Lengkap Formula Matematis Sistem}
\begin{tabularx}{\linewidth}{lXl}
\toprule
\textbf{Nama} & \textbf{Formula} & \textbf{Ref.} \\
\midrule
Kecepatan baru & $v' = \operatorname{clamp}(v + a\Delta t,\; 0,\; v_{\max})$ & \ref{eq:v_update} \\[0.3em]
Posisi baru & $x' = x + v'\cos\theta\,\Delta t$ & \ref{eq:x_update} \\[0.3em]
Lane trajectory & $y_w(t) = y_s + (y_t - y_s)\!\left(\tfrac{1}{2}-\tfrac{1}{2}\cos\pi p\right)$ & \ref{eq:lane_pos} \\[0.3em]
Desired gap & $d_{\text{des}} = s_0 + h\, v_f$ & \ref{eq:d_desired} \\[0.3em]
Spacing error & $e = (x_p - x_f) - d_{\text{des}}$ & \ref{eq:spacing_err} \\[0.3em]
Rel. speed & $v_{\text{rel}} = v_p - v_f$ & \ref{eq:v_rel} \\[0.3em]
CACC accel & $a_{\text{CACC}} = 0.45\,e + 0.35\,v_{\text{rel}} + 0.6\,a_{\text{ff}}$ & \ref{eq:cacc} \\[0.3em]
ACC headway & $h_{\text{safe}} = \max(2.0,\; 1.6\,h)$ & \ref{eq:acc_h} \\[0.3em]
ACC accel & $a_{\text{ACC}} = 0.32\,e + 0.072\,v_{\text{rel}}$ & \ref{eq:acc_a} \\[0.3em]
Latency model & $t_{\text{deliver}} = t_{\text{now}} + L_{\text{ms}}$ & \ref{eq:latency} \\[0.3em]
E2E delay & $T_{\text{E2E}} = L_{\text{ms}} \times (N-1)$ & \ref{eq:e2e} \\[0.3em]
Utilisasi & $U = \min\!\left(100,\frac{N_v \times 0.05}{B_{\text{MHz}}}\times 100\right)$ & \ref{eq:util} \\[0.3em]
SSI & $SSI = \max\!\left(0,\; 1 - |e|/20\right)$ & \ref{eq:ssi} \\[0.3em]
Collision detect & $|x_A{-}x_B| < 4.5\,\wedge\,|y_A{-}y_B| < 2.0$ (meter) & \ref{eq:collision} \\
\bottomrule
\end{tabularx}
\end{table}

% ============================================================
\section{Kesimpulan}
% ============================================================

Penelitian ini berhasil merancang dan mengimplementasikan \textbf{Platform Simulasi 2D
Interaktif Client-Server} yang terintegrasi dengan emulator jaringan 5G URLLC dan
algoritma kontrol CACC. Temuan utama:

\begin{enumerate}
  \item \textbf{Validasi CACC}: Algoritma PD + Feedforward mempertahankan formasi dengan
        SSI mendekati 1.0 pada kondisi ideal ($L < 20$ ms, PLR $< 5\%$).

  \item \textbf{Pengaruh Jaringan}: Peningkatan latensi dan packet loss langsung menurunkan
        SSI. Pada PLR $\geq 15\%$, sistem otomatis beralih ke ACC dengan headway lebih lebar.

  \item \textbf{Efektivitas Dynamic Path Loss}: Model 3GPP berhasil merepresentasikan kondisi
        jaringan jalan raya nyata yang berfluktuasi di antara dua tiang RSU.

  \item \textbf{Superioritas Hybrid Topology}: Topologi Hybrid memberikan kestabilan terbaik
        dibandingkan PF maupun L2A pada skenario gangguan jaringan.

  \item \textbf{Platform Aman}: Simulator memungkinkan pengujian skenario ekstrem tanpa
        risiko fisik dan dengan biaya minimal dibandingkan uji coba nyata.
\end{enumerate}

% ============================================================
\section*{Daftar Pustaka}
\addcontentsline{toc}{section}{Daftar Pustaka}
% ============================================================

\begin{enumerate}[label={[\arabic*]}]
  \item Y.~Zhang, E.~Li, X.~Zhang, and Y.~Zhao, ``A Survey on Vehicular Platooning
        Communication, Control, and Applications,'' \textit{IEEE Trans. Intell. Transp.
        Syst.}, vol.~22, no.~6, pp.~3311--3326, 2021.

  \item J.~Wang and H.~Liang, ``Performance Analysis of V2V Communication for Vehicle
        Platooning Applications,'' \textit{IEEE Access}, vol.~8, pp.~44872--44884, 2020.

  \item 3GPP TS 38.211, ``NR; Physical Channels and Modulation (Release 16),''
        3rd Generation Partnership Project, 2020.

  \item S.~Tsugawa, S.~Jeschke, and S.~E.~Shladover, ``A Review of Truck Platooning
        Projects for Energy Savings,'' \textit{IEEE Trans. Intell. Veh.}, vol.~1, no.~1,
        pp.~68--77, 2016.

  \item V.~Milanés and S.~E.~Shladover, ``Modeling Cooperative and Autonomous Adaptive
        Cruise Control Dynamic Responses Using Experimental Data,''
        \textit{Transp. Res. Part C: Emerg. Technol.}, vol.~48, pp.~285--300, 2014.
\end{enumerate}

\vfill
\begin{center}
  \small\color{gray}\rule{0.5\linewidth}{0.4pt}\\[0.3em]
  Rangkuman Tugas Akhir -- Simulator Platooning V2V Berbasis 5G\\
  \texttt{https://v2v-platooning-simulator.cloud} $\cdot$ Telkom University $\cdot$ 2026
\end{center}

\end{document}
